\chapter{Methodology}

\section{Dataset and Preprocessing}
\label{sec:dataset}
The experimental evaluation in this thesis utilizes the "SmartMeter Energy Consumption Data in London Households" dataset, provided by UK Power Networks as part of the Low Carbon London project \cite{londonSmartMeter2014}. The original dataset covers the period from November 2011 to February 2014 and contains approximately 167 million records from 5,567 representative households in the Greater London area. The consumption readings were recorded at a half-hourly granularity, providing a high-frequency energy profile that is essential for the analysis of Snapshot $z$-anonymity.

\subsection{Data Selection}
Due to the computational volume of the full dataset, a representative time window was selected for the simulations. The experiments utilize a one-week period from November 5, 2012 (00:00:00) to November 11, 2012 (23:30:00). Given the half-hour sampling interval, this corresponds to 336 timestamps within the selected dataset window. This subset allows for the analysis of both weekday and weekend consumption patterns while maintaining a manageable data volume for iterative testing.
The selected subset contains 1,854,193 tuples. It includes readings from 5,529 unique households. Although the project lists 5,567 total participants, the discrepancy of 38 individuals is attributed to missing reports or device downtime during this specific week. As this represents less than 0.7\% of the data population, it does not statistically impact the validity of the privacy analysis.

\subsection{Preprocessing and Schema}
Raw data contained tariff information (such as the stdorToU column indicating Standard or Dynamic Time-of-Use tariffs). As the focus of this research is on the privacy of consumption values rather than billing mechanisms, this column was removed during preprocessing.
The final data schema used for the experiments is defined as a tuple T=⟨ID,t,v⟩, where:
\begin{itemize}
\item \textbf{ID (LCLid):} A unique anonymized identifier for the smart meter (household).
\item \textbf{t (DateTime):} The timestamp of the reading (30-minute granularity).
\item \textbf{v (KWH/hh):} The energy consumption value in kilowatt-hours per half hour.
\end{itemize}
Prior to the experiments, the dataset was checked for missing values. No null entries were found in the selected one-week subset; therefore, no data cleaning regarding missing values was required, and the full sample of 1,854,193 tuples was used for the simulations.

\section{Implementation of Experimental Scenarios}

To verify the hypotheses, the simulation framework was developed using the \textbf{Python 3.13} programming language. The complete source code, including the experiment classes and the evaluation scripts, is publicly available on GitHub\footnote{\url{https://github.com/kaczynskimat/bachelorarbeit}}. The implementation leverages the scientific Python ecosystem to handle the large volume of smart meter data efficiently. Specifically, the \textbf{Pandas} library is utilized for high-performance data manipulation, employing vectorized operations to process the 1.8 million tuples without the performance overhead of iterative loops. Visualization of the resulting data distributions and metrics is handled by \textbf{Matplotlib} and \textbf{Seaborn}.

The codebase follows a modular object-oriented design pattern to ensure reproducibility. Rather than a monolithic script, each experimental scenario is encapsulated in a distinct, standalone Python class. While these classes do not share a common inheritance parent, they adhere to a standardized structural contract, implementing identical methods for data loading, processing, and metric calculation. This design ensures that all strategies are evaluated against the same criteria and produce a standardized dictionary of results. The simulations are orchestrated by a central runner script (\texttt{main.py}), which iterates through the parameter ranges and consolidates the outputs into a single CSV file for comparative analysis.

\subsection{Scenario A: Baseline (Cloud-Only)}
The first scenario represents the centralized snapshot z-anonymity model with a trusted third party as described in~\cite{zanonDataStreams}. It serves as a baseline/control group to compare our remaining results and the effects of the preprocessing steps. In this configuration, no preprocessing or data transformation occurs at either the Edge (Smart Meter) or the Fog (Gateway) layers. Implemented in the \texttt{BaselineExperiment} class, this simulation models Smart Meters transmitting raw, high-precision energy readings (three decimal places) directly upon generation. The intermediate Gateways at the Fog layer perform no filtering or local analysis, acting strictly as passive relays that forward the complete data stream to the Central Entity. Consequently, the snapshot $z$-anonymity algorithm is executed entirely at the Cloud layer. This scenario establishes the benchmark metrics for publication ratio and bandwidth consumption. By identifying how much data is suppressed when raw, high-entropy values from the entire population are subjected to strict anonymity constraints, it provides a "worst-case" reference point against which the effectiveness of distributed preprocessing is measured.



\subsection{Scenario B: Data Generalization (Edge)}
In the second scenario, the privacy-preserving workload is shifted to the Edge layer, simulating a model where the Smart Meter sanitizes data before transmission. The primary objective is to increase the probability of satisfying the $z$-anonymity threshold by reducing the granularity of the consumption values. When the raw readings are less precise, the likelihood of multiple users reporting identical values at the same timestamp increases, thereby reducing data suppression.

This logic is implemented in the \texttt{GeneralizationExperiment} class. Before the privacy check, a scalar rounding transformation is applied to every energy reading $v$. To ensure that the data is mapped to its mathematically nearest representative value, the implementation uses the standard round-half-up formula:

\begin{equation}
    v' = \frac{\lfloor v \cdot 10^p + 0.5 \rfloor}{10^p}
\end{equation}

In this equation, $p$ represents the target precision level (the number of decimal places). The inclusion of the $+0.5$ offset within the floor function $\lfloor \dots \rfloor$ is a deliberate architectural choice to perform rounding rather than simple truncation. Without this offset, the data would always be rounded down, leading to a significant systematic bias and higher information loss. By rounding to the nearest neighbor, the algorithm minimizes the error introduced during generalization, preserving as much utility as possible for the utility provider.

The simulation systematically evaluates three levels of generalization ($p \in \{0, 1, 2\}$) and compares them against the raw dataset resolution ($p=3$). For $p=0$, values are rounded to the nearest integer (e.g., $0.158$ kWh becomes $0.0$ kWh), whereas $p=2$ retains a higher degree of granularity. The resulting quantization error and its impact on data utility are further analyzed using both the standard Normalized Certainty Penalty (NCP) and the refined Effective NCP metric, both of which are described in detail in Section \ref{sec:evaluation_metrics}.

\subsection{Scenario C: Temporal Aggregation (Edge)}
The third scenario evaluates the impact of reducing the temporal granularity of the data stream. Implemented in the \texttt{TemporalAggregationExperiment} class, this strategy simulates an Edge node that buffers multiple readings and transmits a single representative average over a fixed time window $W$. This approach contrasts with the baseline 30-minute reporting interval and is specifically designed to evaluate the trade-off between real-time monitoring and network efficiency.

The implementation utilizes the Pandas \texttt{resample} function to apply non-overlapping "tumbling" windows. For a set of $n$ readings $\{v_1, v_2, \dots, v_n\}$ captured within a window $W$, the Edge node calculates and transmits the arithmetic mean $\bar{v}$:

\begin{equation}
    \bar{v} = \frac{1}{n} \sum_{i=1}^{n} v_i
\end{equation}

The simulation evaluates window sizes of $W \in \{1H, 2H, 4H\}$, corresponding to $n \in \{2, 4, 8\}$ readings per transmission. This preprocessing step has a dual impact on the dataset:
\begin{enumerate}
    \item \textbf{Data Volume:} It directly reduces the total row count of the transmitted dataset by a factor of $n$. This leads to significant \textit{Bandwidth Savings} (as defined in Section \ref{sec:evaluation_metrics}), which is a critical performance objective in resource-constrained \ac{IoT} environments.
    \item \textbf{Value Transformation:} The process of averaging values over a time window aims to reduce the distinctness of usage peaks, theoretically making individual profiles less unique and more similar to the broader population. This research investigates whether this smoothing increases the publication ratio by making users appear more similar, or if the generation of new, high-precision arithmetic means actually increases the sparsity of the dataset, thereby harming data availability.

\end{enumerate}

Unlike the Generalization scenario, which only reduces the precision of a single reading, Temporal Aggregation fundamentally alters the temporal resolution of the data, masking fine-grained behavioral patterns while providing a significant performance boost for the \ac{IoT} architecture.



\subsection{Scenario D: Local Prefiltering (Fog)}
The final scenario evaluates a distributed privacy-preserving architecture that leverages the intermediate Fog layer. This logic is implemented in the \texttt{LocalPrefilteringWithGeneralization} class and represents a multi-stage hybrid approach to data protection.



% To simulate a realistic network topology, the user population ($N=5,529$) is logically partitioned into 56 neighborhoods. To ensure statistical consistency across the majority of the simulation, 55 of these neighborhoods are modeled as standard clusters of 100 households each. The final neighborhood consists of the remaining 29 households. This configuration allows the research to evaluate the preprocessing strategies across a near-uniform distribution while maintaining the integrity of the full real-world dataset. The processing pipeline consists of three distinct stages:


As defined in the system model (Section \ref{sec:system_model}), the user population is logically partitioned into 56 neighborhoods, each monitored by a specific Fog Gateway. The processing pipeline for this localized scenario consists of three distinct stages:



\begin{enumerate}
    \item \textbf{Edge Level Generalization:} Prior to transmission to the gateway, each Smart Meter rounds its energy reading $v$ to a fixed precision of $p=2$ decimal places. This specific precision level was determined based on a preliminary analysis of the dataset. When local prefiltering was applied to the raw, high-precision data ($p=3$), the resulting suppression rates at the Fog layer were untenably high. Only 12\% of tuples survived even at a minimal threshold of $z_{local}=2$, and less than 0.5\% survived at $z_{local}=4$. Because each Gateway only observes a restricted subset of the population (e.g., 100 households), the raw data lacks the necessary density to form local "crowds." Therefore, generalizing the data to two decimal places was chosen as a necessary prerequisite to ensure that a viable volume of data survives the local privacy check.
    \item \textbf{Fog Level Prefiltering:} Each Gateway executes a local snapshot $z$-anonymity check on the generalized data received from its neighborhood. Following the surplus publication model, the Gateway identifies the frequency of each value $v$ within its local subset $U_i$. A tuple is forwarded to the Central Entity only if it satisfies $Count(v, U_i) \geq z_{local}$. The number of tuples forwarded for a given value, $N_{fwd}$, is calculated as:
    \begin{equation}
        N_{fwd}(v, U_i) = \max(0, Count(v, U_i) - (z_{local} - 1))
    \end{equation}
    \item \textbf{Global Anonymity Check:} The Central Entity collects the surplus tuples from all gateways and performs the final global snapshot $z$-anonymity check (as defined in Section~\ref{sec:zanonsnapshot}) before the data is considered fully anonymized and ready for publication.
\end{enumerate}

This scenario evaluates the hypothesis that performing a preliminary privacy check at the Fog layer can significantly reduce the volume of "unique outliers" transmitted over the core network. By suppressing rare values close to the source, the architecture aims to maximize \textit{Bandwidth Savings} while maintaining high global data availability for the Central Entity.


\section{Evaluation Metrics}
\label{sec:evaluation_metrics}

To quantify the trade-offs among privacy, utility, and performance across different architectural scenarios, three primary metrics are used.

\subsection{Publication Ratio}
The Publication Ratio ($PR$) measures the availability of data after the snapshot $z$-anonymity constraints have been applied. It is defined as the percentage of tuples that satisfy the privacy threshold and are released to the end-user or application:
\begin{equation}
    PR = \frac{N_{published}}{N_{total}} \times 100\%
\end{equation}
where $N_{published}$ is the count of tuples satisfying $Count(v, P_t) \geq z$, and $N_{total}$ is the total number of tuples in the original dataset.

\subsection{Bandwidth Savings}
To evaluate the performance benefits of the Fog and Edge layers, we measure the reduction in data transmission volume. Bandwidth Savings ($BS$) is calculated as the percentage of messages suppressed or aggregated before reaching the Central Entity:
\begin{equation}
    BS = \left( 1 - \frac{N_{transmitted}}{N_{total}} \right) \times 100\%
\end{equation}
In the baseline, $BS$ is $0\%$, whereas in temporal aggregation and local prefiltering, this value represents the network load reduction.

\subsection{Information Loss and Effective NCP}
Information loss is measured using the Normalized Certainty Penalty (NCP). The standard NCP normalizes the width of a generalization interval $[L, U]$ against the global range of the attribute:
\begin{equation}
    NCP_{std} = \frac{U - L}{max(v) - min(v)} \times 100\%
\end{equation}

While traditional applications of NCP aggregate the penalties of individually varying interval sizes across all tuples, the generalization technique applied in this thesis (scalar rounding) employs uniformly sized intervals. Specifically, rounding to $p$ decimal places dictates a constant interval width of $U - L = 10^{-p}$ for all records. Consequently, the average information loss across the entire dataset can be deterministically calculated as a single global metric:

\begin{equation}
    NCP_{std} = \frac{10^{-p}}{\max(v) - \min(v)} \times 100\%
    \label{eq:ncp_std_uniform}
\end{equation}


However, smart meter data is characterized by high skewness and extreme outliers. To prevent these rare high-consumption values from artificially inflating the denominator and "diluting" the perceived information loss, this research introduces a refined \textbf{Effective NCP ($NCP_{eff}$)}. This metric evaluates utility loss relative to the range of the typical population. The "Effective Range" is determined through a robust statistical process:
\begin{enumerate}
    \item The median consumption $\tilde{v}$ of the dataset is calculated.
    \item For every reading $v_i$, the absolute deviation from the median is computed: $|v_i - \tilde{v}|$.
    \item A cutoff distance $D$ is identified as the 95th percentile of these deviations.
    % \item The dataset is filtered to include only the 95\% of observations where $|v_i - \tilde{v}| \leq D$.
    \item The dataset is filtered to retain all observations that satisfy the condition:
    \begin{equation}
        |v_i - \tilde{v}| \leq D
    \end{equation}
\end{enumerate}
% The $NCP_{eff}$ uses the range of this filtered subset as the denominator. This ensures that the metric accurately reflects the impact of rounding on the vast majority of households, providing a more rigorous assessment of data utility than the standard global normalization.
The $NCP_{eff}$ uses the range of this filtered subset as its denominator. It is crucial to note a conceptual departure from the standard NCP defined by Xu et al. \cite{UtilityBasedAnonymizationNCP}. This metric measures something fundamentally different than the standard NCP: while $NCP_{std}$ measures information loss relative to the global range, $NCP_{eff}$ measures the \textbf{relative interval width compared to the data span and variance of the typical population} (the central 95\% of users).

% In standard hierarchical generalization, intervals are strictly bounded subsets of the global domain, ensuring $NCP_{std} \leq 100\%$. However, this research applies data quantization (rounding) independent of the data distribution. Consequently, if the forced rounding interval (e.g., 1.0 kWh for integer rounding) exceeds the actual consumption span of the 95\% typical population, the $NCP_{eff}$ will mathematically exceed 100\%. 

In this context, values of $NCP_{eff} >100\%$ are mathematically allowed and do not indicate an error. Instead, they provide a clear signal that the chosen rounding interval is larger than the actual variation of the data within that group. When the metric exceeds 100\%, it indicates that the preprocessing removes too much detail from the data, causing originally distinct household consumption profiles to become indistinguishable. Consequently, the $NCP_{eff}$ serves as an additional indicator of whether a specific rounding precision is appropiate for the typical population within the dataset.
